\chapter*{उत्तर-कुंजी (Answer Key)}
\addcontentsline{toc}{chapter}{उत्तर-कुंजी (Answer Key)}
\markboth{उत्तर-कुंजी}{उत्तर-कुंजी}

यहाँ \textbf{अ (रिक्त स्थान)} एवं \textbf{ब (लघु उत्तरीय)} के उत्तर दिए गए हैं। \textbf{स}
भाग के अधिकांश प्रश्न गतिविधि एवं विचार आधारित हैं, अतः उनके लिए केवल \emph{संकेत/आदर्श
बिंदु} दिए गए हैं - विद्यार्थियों के उत्तर भिन्न हो सकते हैं।

% ---------------- अध्याय 1 ----------------
\section*{अध्याय 1: कृत्रिम बुद्धिमत्ता का परिचय}
\textbf{अ.} 1. Machine Learning (मशीन लर्निंग) \quad 2. Narrow (सँकरी) AI

\noindent\textbf{ब.}
\begin{enumerate}[leftmargin=1.6em]
  \item परम्परागत प्रोग्राम को \emph{नियम + डेटा} दिए जाते हैं और वह नियम लागू करता है;
        ML को \emph{डेटा + उत्तर} दिए जाते हैं और वह नियम स्वयं खोजती है। उदाहरण -
        कैलकुलेटर (परम्परागत) बनाम चेहरा-पहचान (ML)।
  \item विज्ञान में आज डेटा बहुत विशाल है; AI उसे (i) तेज़ी से, (ii) धैर्यपूर्वक तथा
        (iii) बड़े पैमाने (scale) पर विश्लेषित कर ऐसे पैटर्न पकड़ लेती है जो मनुष्य से
        छूट सकते हैं।
\end{enumerate}

\noindent\textbf{स.} \textit{(गतिविधि)} सामान्यतः class बढ़ाने पर, समान मात्रा में
डेटा हो तो accuracy कुछ घट सकती है; अधिक एवं विविध चित्र देने पर सुधरती है।

% ---------------- अध्याय 2 ----------------
\section*{अध्याय 2: भौतिकी में कृत्रिम बुद्धिमत्ता}
\textbf{अ.} 1. रैखिक समाश्रयण / Linear Regression (Least Squares) \quad
2. दूरी/त्रुटि (error) \quad 3. क्रमबद्ध (Systematic) त्रुटि

\noindent\textbf{ब.}
\begin{enumerate}[leftmargin=1.6em]
  \item गैलीलियो ने झुके तल पर गेंद के समय एवं दूरी का \emph{डेटा} लेकर पैटर्न
        ($s \propto t^2$) खोजा - अर्थात् भौतिकी का नियम डेटा में पैटर्न खोजने से ही मिला।
  \item regression डेटा-बिंदुओं के बीच सर्वोत्तम रेखा खोजता है; उसकी ढलान ही नियम का
        स्थिरांक होती है। जैसे बल बनाम विस्तार के डेटा से ढलान = स्प्रिंग नियतांक $k$
        ($F = kx$)।
  \item LHC प्रति सेकंड लगभग एक अरब टकराव एवं भारी डेटा बनाता है; इनमें से दुर्लभ
        महत्वपूर्ण घटनाएँ छाँटना केवल AI से ही सम्भव है।
\end{enumerate}

\noindent\textbf{स.}
\begin{enumerate}[leftmargin=1.6em]
  \item \textit{(गतिविधि)} बल बनाम विस्तार की सर्वोत्तम रेखा की ढलान ही $k$ है।
  \item $m = 2,\; c = 0$; सम्बन्ध $y = 2x$ अर्थात् $y$, $x$ के अनुक्रमानुपाती - जैसे
        हुक का नियम $F = kx$।
  \item \textit{(गतिविधि)} त्वरण-ग्राफ़ प्रायः आवर्ती (periodic) दिखता है; उत्तर
        एकत्र डेटा पर निर्भर।
\end{enumerate}

% ---------------- अध्याय 3 ----------------
\section*{अध्याय 3: रसायन विज्ञान में कृत्रिम बुद्धिमत्ता}
\textbf{अ.} 1. SMILES (स्माइल्स संकेतन) \quad 2. वर्गीकरण (classification) \quad
3. आभासी छँटाई (Virtual Screening)

\noindent\textbf{ब.}
\begin{enumerate}[leftmargin=1.6em]
  \item मेंडेलीफ़ ने ज्ञात तत्वों की \emph{प्रवृत्ति (trend/pattern)} देखकर अनदेखे
        तत्वों के गुणों का पूर्वानुमान किया - ML भी ज्ञात डेटा के पैटर्न से अनदेखे मान
        का पूर्वानुमान करता है।
  \item classification में वस्तु को श्रेणी में बाँटा जाता है (जैसे अम्ल/क्षार);
        prediction में कोई सतत मान आँका जाता है (जैसे क्वथनांक)।
  \item (i) आभासी छँटाई (virtual screening), (ii) गुण/विषाक्तता-पूर्वानुमान, तथा
        (iii) नए अणुओं का अभिकल्प (molecule design)।
\end{enumerate}

\noindent\textbf{स.} \textit{(गतिविधि)} आवर्त में बाएँ से दाएँ त्रिज्या घटती है; 1.
पूर्वानुमान-त्रुटि छोटी आनी चाहिए। 2. MolView से SMILES लिखें (जैसे एथेनॉल = \texttt{CCO})।
3. निबंध - virtual screening एवं गुण-पूर्वानुमान से समय व लागत घटती है (जैसे Halicin)।

% ---------------- अध्याय 4 ----------------
\section*{अध्याय 4: जीव विज्ञान में कृत्रिम बुद्धिमत्ता}
\textbf{अ.} 1. A, T, G, C \quad 2. AlphaFold \quad 3. Bioinformatics (जैव-सूचना विज्ञान)

\noindent\textbf{ब.}
\begin{enumerate}[leftmargin=1.6em]
  \item DNA केवल चार अक्षरों (A, T, G, C) के क्रम से बना एक विशाल अनुक्रम है, जिसे
        कंप्यूटर पढ़ एवं विश्लेषित कर सकता है - इसीलिए यह एक ``डेटा'' है।
  \item AI चित्रों में रोग के सूक्ष्म चिह्न दिखाकर एक \emph{द्वितीय राय (second
        opinion)} देती है, परन्तु अंतिम निर्णय सदैव चिकित्सक का होता है।
  \item AlphaFold ने प्रोटीन-वलन (protein folding) की दशकों पुरानी समस्या हल की,
        जिससे नई दवाओं एवं एंज़ाइमों की खोज बहुत तेज़ हो गई।
\end{enumerate}

\noindent\textbf{स.}
\begin{enumerate}[leftmargin=1.6em]
  \item \textit{(गतिविधि)} अधिक एवं विविध चित्र देने पर सटीकता (accuracy) बढ़ती है।
  \item \texttt{ATGCATGC} में A = 2, T = 2, G = 2, C = 2; \;
        GC content $= \dfrac{2+2}{8}\times 100 = 50\%$।
  \item \textit{(निबंध)} ऐप करोड़ों तस्वीरों पर प्रशिक्षित होकर प्रजाति पहचानते हैं;
        नागरिक डेटा-संग्रह में भाग ले सकते हैं।
\end{enumerate}

% ---------------- अध्याय 5 ----------------
\section*{अध्याय 5: डेटा, प्रयोग एवं AI उपकरण}
\textbf{अ.} 1. garbage \quad 2. सफ़ाई (Cleaning) \quad 3. Google Colab

\noindent\textbf{ब.}
\begin{enumerate}[leftmargin=1.6em]
  \item शुद्धता (accuracy), विविधता (diversity) एवं पर्याप्त मात्रा (sufficient
        quantity)।
  \item संग्रह $\to$ सफ़ाई $\to$ विश्लेषण $\to$ मॉडलन $\to$ निष्कर्ष।
  \item पक्षपात (bias) तब होता है जब डेटा एकपक्षीय हो और परिणाम झुक जाएँ - जैसे केवल
        खेल-दल से मापी गई ``औसत लम्बाई'' पूरी कक्षा के लिए ग़लत होगी।
\end{enumerate}

\noindent\textbf{स.}
\begin{enumerate}[leftmargin=1.6em]
  \item \textit{(गतिविधि)} दिए गए डेटा पर लगभग $m \approx 50.5,\; c \approx 0.01$
        प्राप्त होगा (मान थोड़े भिन्न हो सकते हैं)।
  \item \textit{(गतिविधि)} सामान्यतः धनात्मक सहसम्बन्ध (positive correlation) -
        लम्बाई बढ़ने पर भुजा-विस्तार भी बढ़ता है।
  \item \textit{(रिपोर्ट)} पाँचों चरणों का पालन एवं डेटा की सीमाओं का उल्लेख अपेक्षित।
\end{enumerate}

% ---------------- अध्याय 6 ----------------
\section*{अध्याय 6: AI की नैतिकता एवं भविष्य}
\textbf{अ.} 1. hallucination (मतिभ्रम) \quad 2. सहसम्बन्ध (correlation) \quad
3. पारदर्शिता (Transparency)

\noindent\textbf{ब.}
\begin{enumerate}[leftmargin=1.6em]
  \item AI केवल अपने प्रशिक्षण-डेटा जितना ही ``जानती'' है और कभी-कभी पूरे आत्मविश्वास
        से ग़लत/गढ़ा हुआ (hallucinated) उत्तर देती है; साथ ही वह केवल सहसम्बन्ध देखती
        है, कारण नहीं।
  \item पारदर्शिता (Transparency), निष्पक्षता (Fairness), जवाबदेही (Accountability)
        एवं निजता (Privacy)।
  \item चिकित्सा/जीव विज्ञान में व्यक्तिगत डेटा (जैसे चिकित्सा-रिकॉर्ड, DNA) प्रयोग
        होता है; उसकी सुरक्षा एवं केवल अनुमति से उपयोग वैज्ञानिक का नैतिक उत्तरदायित्व है।
\end{enumerate}

\noindent\textbf{स.} \textit{(गतिविधि/निबंध)} 1. कठिन उदाहरणों (कम रोशनी, उलटी वस्तु)
पर मॉडल की चूक देखें। 2-3. संतुलित दृष्टिकोण अपेक्षित - AI एक उपकरण है, अंतिम निर्णय
एवं नैतिक उत्तरदायित्व मनुष्य का।

% ---------------- बहुविकल्पीय उत्तर ----------------
\section*{बहुविकल्पीय प्रश्नों के उत्तर (MCQ Answers)}
\noindent
\textbf{अध्याय 1:} 1.~(A) \quad 2.~(B) \quad 3.~(B) \quad 4.~(C)\\
\textbf{अध्याय 2:} 1.~(B) \quad 2.~(B) \quad 3.~(A) \quad 4.~(A)\\
\textbf{अध्याय 3:} 1.~(A) \quad 2.~(B) \quad 3.~(A) \quad 4.~(B)\\
\textbf{अध्याय 4:} 1.~(C) \quad 2.~(A) \quad 3.~(A) \quad 4.~(B)\\
\textbf{अध्याय 5:} 1.~(B) \quad 2.~(A) \quad 3.~(A) \quad 4.~(A)\\
\textbf{अध्याय 6:} 1.~(A) \quad 2.~(A) \quad 3.~(C) \quad 4.~(B)
